\documentclass[twocolumn]{aastex631}
\begin{document}
\title{The reionization ionizing-photon-budget shortfall for star-forming galaxies is not robust to systematics at z\~{}5 (jwsthigh systematic corner)}
\author{NebulaMind Lab (autonomous pipeline)}
\affiliation{NebulaMind Lab, \url{https://lab.nebulamind.net}}
\begin{abstract}
Reionization ionizing-photon-budget reconciliation at z\~{}5: star-forming galaxies require f\_esc=0.009 (+0.008/-0.004) to close the budget (Madau-Dickinson SFRD, log xi\_ion=25.5+/-0.15, clumping C=2-5, JWST-SFRD tail), versus indirect-proxy-inferred f\_esc=0.062 (+0.110/-0.039) from LzLCS O32/beta calibrations. Median delta(required-inferred)=-0.051 dex-frac (16-84\%: -0.161 to -0.012); 5\% of the systematic MC shows a shortfall. Conclusion: the budget CLOSES within the systematic, and the sign holds under both O32 and beta calibrations. The result is bounded by the xi\_ion x clumping x proxy-calibration systematic, not by statistics. Generated autonomously from public data (jwst) via the NebulaMind Lab runner: a bounded, reproducible, descriptive study.
\end{abstract}
\section{Introduction}
Recent studies have highlighted a potential crisis in reconciling the ionizing photon budget during reionization, particularly with the advent of new observations from advanced telescopes [Muñoz2024]. This has led to increased scrutiny of the assumptions underlying our understanding of this critical period in cosmic history. The discrepancy between the expected and observed ionizing photon production has sparked discussions on the role of star-forming galaxies and their escape fractions [Park2022, Davies2021].
\section{Data and method}
To address this issue, we employ a literature-anchored budget calculation that does not rely on new survey catalog data. Instead, we utilize established values from published works: the cosmic SFRD is derived from the Madau \& Dickinson (2014) analytic fitting function, while xi\_ion and O32/beta f\_esc proxy calibrations are adopted from LzLCS studies [Chisholm+22, Flury+22; Simmonds+24]. Our method focuses on reconciling the ionizing-photon-budget using these parameters.
\section{Result}
Our analysis reveals that at z\~{}5, star-forming galaxies must have an escape fraction of f\_esc=0.009 (+0.008/-0.004) to balance the reionization photon budget, assuming a Madau-Dickinson SFRD, log xi\_ion=25.5±0.15, and clumping factor C=2-5. This value is lower than the indirect-proxy-inferred f\_esc=0.062 (+0.110/-0.039) obtained from LzLCS O32/beta calibrations. The median difference between required and inferred escape fractions is -0.051 dex-frac, with a 16-84\% range of -0.161 to -0.012. Notably, 5\% of systematic Monte Carlo simulations indicate a shortfall in the budget.
\begin{figure}\centering\includegraphics[width=\columnwidth]{result.png}\caption{The reionization ionizing-photon-budget shortfall for star-forming galaxies is not robust to systematics at z\~{}5 (jwsthigh systematic corner)}\end{figure}
\section{Caveats}
It is essential to acknowledge the limitations of our approach. The reliance on literature values and calibrations introduces uncertainties tied to the assumptions and methodologies of previous studies. Additionally, our analysis does not account for potential variations in xi\_ion or clumping factor across different galaxy populations. Furthermore, the use of a single selection criterion and uncalibrated measurements may introduce biases that affect the accuracy of our results. A more comprehensive understanding will require incorporating diverse data sets and refining calibration techniques to better constrain these critical parameters.
\section*{References}
\footnotesize
[Muoz2024] Muñoz et al. (2024). Reionization after JWST: a photon budget crisis?. bibcode 2024MNRAS.535L..37M \par [Park2022] Park et al. (2022). Calibrating excursion set reionization models to approximately conserve ionizing photons. bibcode 2022MNRAS.517..192P \par [Duncan2015] Duncan et al. (2015). Powering reionization: assessing the galaxy ionizing photon budget at z < 10. bibcode 2015MNRAS.451.2030D \par [Davies2021] Davies et al. (2021). The Predicament of Absorption-dominated Reionization: Increased Demands on Ionizing Sources. bibcode 2021ApJ...918L..35D \par [Madau2017] Madau (2017). Cosmic Reionization after Planck and before JWST: An Analytic Approach. bibcode 2017ApJ...851...50M

\end{document}
